%%%
%%% 22画
%%%
\section*{22画}\addcontentsline{toc}{section}{22画}\addcontentsline{loh}{figure}{\#\#\#\# 22画}

%%%%%%%%%% 蘸 %%%%%%%%%%
\subsection*{蘸}\addcontentsline{loh}{figure}{蘸}

\begin{Entry}{蘸}{22}{⾋}
  \begin{Phonetics}{蘸}{zhan4}[][HSK 7-9]
    \definition{v.}{mergulhar em (tinta, molho, etc.); mergulhar brevemente em líquidos, pós ou pastas e depois retirar}
  \end{Phonetics}
\end{Entry}

%%%%%%%%%% 镶 %%%%%%%%%%
\subsection*{镶}\addcontentsline{loh}{figure}{镶}

\begin{Entry}{镶}{22}{⾦}
  \begin{Phonetics}{镶}{xiang1}[][HSK 7-9]
    \definition{v.}{para incrustar; cravar; montar | marginar; orlar; rodear; decorar; colocar uma moldura | inserir; integrar}
  \end{Phonetics}
\end{Entry}

\begin{Entry}{镶嵌}{22,12}{⾦,⼭}
  \begin{Phonetics}{镶嵌}{xiang1qian4}[][HSK 7-9]
    \definition{v.}{incrustar; cravar; montar; incorporar um objeto em outro}
  \end{Phonetics}
\end{Entry}

%%%%% EOF %%%%%

